\documentclass[12pt]{article}
\usepackage{lingmacros}
\usepackage{tree-dvips}
\usepackage{graphicx}

\begin{document}

% Title page
\begin{titlepage}
\begin{center}

% University logo and header
\vspace*{0.3cm}
\begin{figure}
    \centering
    \includegraphics[width=0.1\linewidth]{uniandeslogo.png}
\end{figure}
\vspace{0.3cm}

{\Large\textbf{UNIVERSIDAD DE LOS ANDES}}
\vspace{0.5cm}

{\large\textbf{FACULTAD DE INGENIERÍA}}

\vspace{0.7cm}

{\normalsize Ingeniería de Sistemas y Computación}

\vspace{0.3cm}

{\normalsize }

\vspace{0.3cm}

% {\normalsize ISIS3510 - Mobile Applications Development}

% \vspace{3cm}

% Report title
{\LARGE\textbf{Proyecto de Grado 2026-20}}

\vspace{3cm}

% App name
{\Large\textbf{QuickR}}
\vspace{0.3cm}

\vspace{1cm}

% Authors section
%\textbf{Submitted by:}

%\vspace{0.5cm}

{\textbf{Mateo López Cuartas}}

\vspace{4cm}


\vfill

% Date and university location
\today

\vspace{0.5cm}

{\normalsize Bogotá D.C., Colombia}

\end{center}
\end{titlepage}

\newpage
\tableofcontents
\newpage
\setlength{\parindent}{20pt}

\section{Introduction}
\subsection{Context}
Ever since the early 2000s, the internet has had a significant impact on global communication technologies. By the year 2000, 
51\% of global telecommunications networks were transmitted via the internet, a figure that rose to 97\% by 2007 (Hilbert \& López, 2011).
This reliance on the internet has been steadily increasing over time, culminating in 2020 when quarantine and lockdown measures due to the pandemic 
made the internet an essential part of daily life for most people worldwide. This drastic change in people's lives led to an all-time high demand for 
applications that connect different individuals through digital means, such as Meta's Metaverse, the increased demand for virtual meeting software 
(Zoom, Google Meet, Microsoft Teams), and mass communication services like Discord and Slack.

It is in this context, where in-person activities take a secondary role compared to online activities, that the need arises to adapt these daily activities to a new digital environment; in this case, broadcasting and sharing multimedia content via streaming from one device for others to view on their devices in real time, regardless of their geographical location. This streaming became popular through formal/academic means such as Zoom and Microsoft Teams, while the more casual media streaming needs were dominated by platforms like Twitch and Discord (the latter, although primarily used as a social messaging platform, also offers VoIP and video streaming services, 
These services are the main focus of this paper). However, the rapid expansion of video traffic on the internet has created a strong need to develop robust protocols that work efficiently under suboptimal connectivity conditions, as well as infrastructures that are sufficiently robust and scalable to meet the demand for these transmissions.

Es en este contexto, donde las actividades presenciales toman un rol secundario comparado a actividades en línea, surge la necesidad de adaptar estas actividades que previamente eran cotidianas a un nuevo ambiente digital; en este caso el transmitir contenido multimedia de forma de streaming desde un dispositivo para que otras personas lo vean desde sus dispositivos en tiempo real. Este  streaming se popularizó mediante medios formales/académicos como Zoom y Microsoft Teams, mientras que los medios informales de streaming fueron dominados por plataformas como Twitch y Discord (este último, aunque principalmente usado como plataforma social de mensajería, también ofrece servicios de VoIP y transmisión de video, que son los que se deben enfatizar para este proyecto). Sin embargo, la rápida expansión del tráfico de vídeo en internet ha generado una fuerte necesidad de desarrollar protocolos robustos que funcionen de forma eficiente en condiciones de conectividad subóptimas, así como de infraestructuras lo suficientemente sólidas y escalables para satisfacer la demanda de estas transmisiones. 

Estas necesidades, acompañadas de la cantidad masiva de usuarios a nivel global que usan estos servicios, ha llevado a una necesidad por parte de los proveedores de estos servicios a restringir algunas funcionalidades y características del servicio detrás de Paywalls, como por ejemplo Discord, que no permite hacer una transmisión de más de 720p x 30fps sin pagar la subscripción mensual de Nitro (\%9.99 al mes para Agosto de 2026). Sin embargo, esto genera un contraste interesante respecto de la relación entre oferta y demanda de estos servicios de streaming. La demanda de estos servicios de streaming, exacerbada por la pandemia, es suplida casi exclusivamente por proveedores de servicios con suscripciones y tarifas limitadas, lo cual resulta un tanto desconcertante. Ante este panorama de uso cotidiano e intensivo a nivel global, sería razonable contar con una oferta lo suficientemente amplia para democratizar el acceso libre a estos servicios, sin depender de un único proveedor. Esto implica que, entre más oferta haya, más posibilidades tiene el usuario de elegir por su cuenta y no depender de un único modelo.



\end{document}